% RQ3: applied pacing delay against the delay the same admissibility condition
% requires on measured inputs.  Companion to
% Table~\ref{tab:rq3_pacing_calibration}, which carries the summary statistics
% and the full definition of d*.
%
% Panels (a) and (b) are the calibration view: one point per analyzed
% transition with a positive requirement, x the measured requirement d*, y the
% applied delay d, and the dashed diagonal exact calibration.  Above the
% diagonal the controller paced more than the trace required, below it less.
% The +/-100 ms guides are drawn because a transition inside that band is
% calibrated to within roughly one frame interval, which is the granularity at
% which a delay is actuated at all.
%
% ---------------------------------------------------------------------------
% Why the d* = 0 transitions are not plotted here
% ---------------------------------------------------------------------------
% An earlier revision drew them as a hollow column on the vertical axis -- 277
% points in (a) and 278 in (b), 75.3% and 79.9% of the paced transitions -- on
% the reasoning that a reader who cannot see the column cannot judge how much
% of the picture it is.  That reasoning was right about disclosure and wrong
% about where the disclosure belongs.  Two things settle it:
%
%   1. Those points carry no calibration information.  The grammar of these
%      panels is distance from the diagonal, and a point at x = 0 has no such
%      distance; all it shows is the marginal distribution of d.  Rendered at
%      277 overplotted marks on one x it also saturates into a solid bar, so
%      even that marginal cannot be read off it.
%   2. Panel (c) already contains every one of them, quantitatively.  For a
%      transition with d* = 0 the error d - d* is exactly d, so the column is
%      the part of (c) between the step at zero and the right tail.  Removing
%      it from (a) and (b) routes each quantity to the panel whose grammar fits
%      it; it does not remove anything from the figure.
%
% The restriction is therefore declared on the x label, where it cannot be
% mistaken for a hidden filter, and the count plus the pointer to (c) are in
% the caption.  What the column must NOT be labelled, wherever it appears, is
% "unnecessary pacing".  d* is computed on the realized trace, which pacing
% already shaped; removing those delays would deepen the backlog and lengthen
% later Drafts, so the requirement of the unpaced trajectory is at or above d*.
% RQ1(b) supplies the counterfactual directly: with pacing removed the same
% conditions time out.
%
% The caption now makes that case with a measurement rather than with the
% closed-loop argument alone.  On the same eligible set, every one of the 277
% and 278 zero-requirement paced transitions had a positive measured backlog
% (RQ3Pacing.realBacklogMs), a measured time-to-free
% (PacingReplay.realQueueWaitMs) with P10/P50/P90 of 1,925/3,197/4,006 ms in
% 12MP and 1,589/3,069/4,258 ms in 24MP, and realQueueDepth >= 1 on 100% and
% 99.6% of them (median 4 waiting Drafts).  The applied delay covered a median
% 11% and 7% of that measured wait, and stayed below it on 100% and 96.8%.
% d* = 0 therefore does not mean the pipeline was idle; it means the newest
% committed capture's deadline was not yet threatened on a trace pacing had
% already flattened.  Recomputed from the same workbooks and the same run and
% transition eligibility as Table~\ref{tab:rq3_pacing_calibration}; the 12MP
% figures reproduce that table's 1,167 analyzed / 368 paced / 94 required
% exactly, and the 24MP figures come from 1,036 analyzed before the two
% watchdog-truncated transitions are dropped, which do not fall in this subset
% (278 either way).
%
% Positive-requirement points are split by whether admission demoted the
% sequence after the pacing decision was taken, because that is the first
% objection to the surplus: a delay priced for M+S looks over-large once
% admission has removed M from the Draft that ran.  The split answers it with
% data rather than argument, and the answer is that the effect does not run one
% way.  Demoted transitions are 14 of 94 in (a) and 22 of 73 in (b); their
% coverage is 78.6% against 95.0% for undemoted in (a) but 86.4% against 84.3%
% in (b), and their median surplus among paced transitions is higher in (a)
% (416 vs 282 ms) and lower in (b) (243 vs 274 ms).  Coordination therefore does
% not systematically explain the surplus, which is what leaves the
% Draft-duration reserve as the explanation the table reports.  The subsets are
% small, so the split is presented as a marker class to be inspected rather than
% as a table column with statistics.
%
% Panel (c) is the same quantity as a distribution over every analyzed
% transition rather than only the plotted ones, which is what makes the tail
% readable.  The step at zero (68.2% in 12MP, 66.1% in 24MP) is the transitions
% where neither the controller nor the measured requirement asked for a delay;
% the mass to the left of zero is under-pacing and is small (0.6% and 1.1%); the
% right tail is the surplus that Table~\ref{tab:rq3_pacing_calibration}
% attributes to the Draft-duration reserve.
%
% Axis range.  Both scatter panels use an identical 0-800 by 0-1500 ms box so
% the two conditions can be compared by eye.  The box is no longer square
% because the two axes no longer span the same data: no requirement in either
% panel exceeds 727 ms, while applied delays reach 2010 ms, so a square window
% wide enough for the y data would leave more than half the x range empty --
% which is what the previous 0-1500 square did.  The diagonal is consequently
% not at 45 degrees; it is still the exact-calibration locus, which is what the
% panel is read against.  Two of the 73 plotted points in (b) exceed 1500 ms on
% the y axis (largest 2010 ms) and are clipped; no point in (a) is clipped.
% Panel (c) keeps the 1500 ms window on the error axis, where 0.58% of the 24MP
% transitions (largest error 1909 ms) fall outside it and none of the 12MP ones
% do.
%
% Point counts: (a) 94 with a positive requirement (80 undemoted, 14 demoted)
% of 1,167 analyzed transitions; the other 1,073 are the 277 with d* = 0 and
% the 796 at the origin, and neither group is drawn here.  (b) 73 (51 and 22)
% of 1,034, with 278 and 683.
%
% Data: data/rq3/calibration/, generated from
% SM-S948U_metrics_{12MP_normal,24MP_memory}_0803.xlsx (ML commit 99aae0a).
% Run and transition eligibility is documented in
% tables/tab_rq3_pacing_calibration.tex and is identical here.

\begin{figure*}[t]
  \centering
  \providecommand{\rqThreeCalDir}{data/rq3/calibration}

  \tikzset{
    rq3cal required/.style={
      mark=*, mark size=1.25pt, only marks,
      draw=black, fill=black, fill opacity=0.55, draw opacity=0.75},
    rq3cal demoted/.style={
      mark=triangle, mark size=1.9pt, only marks,
      draw=blue!70!black, fill=white, line width=0.5pt, draw opacity=0.9},
    rq3cal diag/.style={black, densely dashed, line width=0.6pt},
    rq3cal band/.style={gray!60, dotted, line width=0.45pt},
    rq3cal ecdf a/.style={black, solid, line width=1.0pt},
    rq3cal ecdf b/.style={black!55, densely dashed, line width=1.0pt},
    rq3cal zeroline/.style={gray!70, solid, line width=0.4pt},
  }

  % #1: file stem of the condition.
  \newcommand{\rqthreecalscatter}[1]{%
    \addplot[rq3cal band,forget plot,domain=0:800,samples=2] {x+100};
    \addplot[rq3cal band,forget plot,domain=100:800,samples=2] {x-100};
    \addplot[rq3cal diag,forget plot,domain=0:800,samples=2] {x};
    \addplot[rq3cal required,forget plot]
      table[x=required_ms,y=applied_ms,col sep=comma]
      {\rqThreeCalDir/scatter_#1_req_kept.csv};
    \addplot[rq3cal demoted,forget plot]
      table[x=required_ms,y=applied_ms,col sep=comma]
      {\rqThreeCalDir/scatter_#1_req_changed.csv};}

  \newsavebox{\rqthreecalfigbox}
  \begin{lrbox}{\rqthreecalfigbox}
  \begin{tikzpicture}
    \begin{groupplot}[
      % Panels (a) and (c) both carry a rotated y label, so the group separation
      % has to clear a label plus a tick column rather than a tick column alone.
      group style={group size=3 by 1, horizontal sep=0.075\textwidth},
      % scale only axis keeps the three plot boxes identical regardless of how
      % wide each one's tick labels are, so the two scatter panels stay square
      % and therefore comparable.
      scale only axis=true,
      width=0.215\textwidth,
      height=0.215\textwidth,
      tick label style={font=\tiny},
      label style={font=\scriptsize},
      title style={font=\scriptsize},
      grid=major,
      major grid style={gray!20,line width=0.2pt},
      axis line style={line width=0.35pt},
      tick style={line width=0.35pt},
      tick align=outside,
      tick pos=left,
      enlargelimits=false,
      clip=true,
    ]

      % The restriction to d* > 0 is stated on the x label rather than in the
      % caption alone: a filter that is visible on the axis it applies to
      % cannot be read as a hidden one.
      \nextgroupplot[
        title={(a) 12MP normal},
        xlabel={Required delay \(d^{*}>0\) (ms)\\
                {\tiny 277 with \(d^{*}=0\) omitted --- see (c)}},
        xlabel style={align=center},
        ylabel={Applied delay \(d\) (ms)},
        xmin=0,xmax=800,ymin=0,ymax=1500,
        xtick={0,200,400,600,800},ytick={0,500,1000,1500},
      ]
      \rqthreecalscatter{12mp_normal}

      \nextgroupplot[
        title={(b) 24MP memory pressure},
        xlabel={Required delay \(d^{*}>0\) (ms)\\
                {\tiny 278 with \(d^{*}=0\) omitted --- see (c)}},
        xlabel style={align=center},
        xmin=0,xmax=800,ymin=0,ymax=1500,
        xtick={0,200,400,600,800},ytick={0,500,1000,1500},
      ]
      \rqthreecalscatter{24mp_memory}

      \nextgroupplot[
        title={(c) Calibration error},
        xlabel={\(d-d^{*}\) (ms)},
        ylabel={Cumulative fraction},
        xmin=-600,xmax=1500,ymin=0,ymax=1.02,
        xtick={-500,0,500,1000,1500},
        ytick={0,0.25,0.5,0.75,1},
        yticklabel style={/pgf/number format/fixed},
      ]
      \addplot[rq3cal zeroline,forget plot] coordinates {(0,0) (0,1.02)};
      \addplot[rq3cal ecdf a,const plot,no marks,forget plot]
        table[x=error_ms,y=cdf,col sep=comma]
        {\rqThreeCalDir/error_ecdf_12mp_normal.csv};
      \addplot[rq3cal ecdf b,const plot,no marks,forget plot]
        table[x=error_ms,y=cdf,col sep=comma]
        {\rqThreeCalDir/error_ecdf_24mp_memory.csv};
    \end{groupplot}

    % Shared legend: the two marker classes explain panels (a) and (b), the two
    % line styles explain panel (c).
    \coordinate (rqthreecallegend) at
      ($(group c1r1.south)!0.5!(group c3r1.south)$);
    \matrix[
      matrix of nodes,
      nodes={font=\scriptsize,anchor=west},
      column sep=2pt,
      row sep=0pt,
      anchor=north,
    % -7.4ex rather than -5.2ex: the scatter x labels now carry a second line
    % declaring the d* = 0 omission, and the legend has to clear it.
    ] at ([yshift=-7.4ex]rqthreecallegend) {
      \draw[draw=black,fill=black,fill opacity=0.55]
        (0,0) circle (1.25pt); & \(d^{*}>0\) &
      \draw[draw=blue!70!black,fill=white,line width=0.5pt]
        (90:1.9pt)--(210:1.9pt)--(330:1.9pt)--cycle;
        & \(d^{*}>0\), demoted &
      \draw[rq3cal ecdf a] (0,0)--(0.38,0); & 12MP normal &
      \draw[rq3cal ecdf b] (0,0)--(0.38,0); & 24MP mem.\ pressure \\
    };
  \end{tikzpicture}
  \end{lrbox}
  \resizebox{\textwidth}{!}{\usebox{\rqthreecalfigbox}}

  \caption{RQ3 pacing calibration. Panels (a) and (b) plot the \(94\) and
  \(73\) analyzed transitions whose measured requirement is positive, one point
  each; the dashed diagonal is exact calibration and the dotted guides mark
  \(\pm100\)\,ms. Triangles mark transitions whose sequence admission demoted
  after the pacing decision, so that the coordination effect on the surplus is
  visible; it does not run one way, being adverse in (a) and neutral in (b). A
  further \(277\) and \(278\) paced transitions have \(d^{*}=0\) on the realized
  trace and are not plotted here; because their error \(d-d^{*}\) is exactly
  \(d\), they are the part of panel~(c) between the step at zero and the right
  tail. They are not idle pacing: every one of them had a positive measured
  backlog, the median measured time until the Draft pipeline was next free was
  \(3.2\) and \(3.1\)\,s, and the applied delay covered a median \(11\%\) and
  \(7\%\) of that wait. Because the trace was already shaped by pacing, the
  requirement of the unpaced trajectory is at or above \(d^{*}\), and
  Table~\ref{tab:rq1_ablation} shows that removing pacing reintroduces Capture
  Timeouts under the same conditions. Panel~(c) gives the calibration error over
  all \(1{,}167\) and \(1{,}034\) analyzed transitions, the transitions at the
  origin included: \(0.6\%\) and \(1.1\%\) of them lie left of zero, and the
  right tail is the surplus that Table~\ref{tab:rq3_pacing_calibration}
  attributes to the Draft-duration reserve.}
  \label{fig:rq3_pacing_calibration}
\end{figure*}
